\subsection{Proof of Surjectivity}
\label{subsections:techniques:universal:surjective}



\paragraph{Description.}
A proof of surjectivity proves that a function is onto, %
meaning that every element of the codomain is an output. %
Precisely, it proves that a function $f : X \to Y$ satisfies %
$\forall y \in Y$, $\exists x \in X$ such that $f(x) = y$. %
Generally, a proof of surjectivity mirrors the following format:
\begin{enumerate}{
    \item \textbf{Setup.} %
        Let $y \in Y$.
    \item \textbf{Verification.} %
        Define $x$ in terms of $y$ such that $x \in X$, %
        and show that $f(x) = y$.
    \item \textbf{Conclusion.} %
        Conclude the proof, optionally by summarizing what has been found.
}
\end{enumerate}



\paragraph{Advice.}
\begin{itemize}{
    \item Use this proof for a statement that a function is onto, %
        also called surjective.
    \item Solve $f(x) = y$ for $x$ to find $x$ in terms of $y$, %
        then confirm that $x \in X$ for every $y \in Y$.
}
\end{itemize}



\paragraph{Example.}
\begin{exercise}{
    Let $f : \R \to \R$ be defined by $f(x) = 2x + 1$. %
    Prove that $f$ is onto.
}
\end{exercise}
\begin{proof}{
    This statement can be proven by showing that $f$ is onto.

    \begin{annotatedpanel}[Setup.]{
        Let $y \in \R$.
    }{
        \textbf{Step 1.} Let $y \in Y$. %
        Here, $Y = \R$.
    }
    \end{annotatedpanel}

    \begin{annotatedpanel}[Verification.]{
        Let $x \in \R$ such that $x = \frac{y - 1}{2}$. %
        Then $f(x) = 2\left(\frac{y - 1}{2}\right) + 1 = y$.
    }{
        \textbf{Step 2.} Define $x$ in terms of $y$ such that $x \in X$, %
        and show that $f(x) = y$. %
        Here, solving $2x + 1 = y$ gives $x = \frac{y - 1}{2}$, %
        which is in $\R$ for every $y \in \R$.
    }
    \end{annotatedpanel}

    \begin{annotatedpanel}[Conclusion.]{
        Thus, $\forall y \in \R$, $\exists x \in \R$ such that $f(x) = y$. %
        Therefore, the statement holds by the definition of onto. \qedhere
    }{
        \textbf{Step 3.} Conclude the proof. %
        Here, the summary states what has been found: %
        every element of $\R$ is an output of $f$.
    }
    \end{annotatedpanel}
}
\end{proof}